\section{Integral Domains, Fields of Fractions, Prime Ideals}

Notice that in the ring \(\flatfrac{\ZZ}{6}\), we have \(2 \cdot 3 = 0\) but \(2 \neq 0\) and \(3 \neq 0\).

\begin{definition}[Integral Domain, Zero Divisor]
    \label{def:integral-domain-zero-divisor}%
    A non-zero ring \(R\) is an \emph{integral domain}, if for all \(a, b \in R\), we have
    \[
        a b = 0 \implies a = 0 \qor b = 0.
    \]

    A \emph{zero divisor} if a non-zero \(a\), such that there exists \(b \neq 0\) if \(ab = 0\).
\end{definition}

For example, in \(\flatfrac{\ZZ}{6}\), we say both \(2\) and \(3\) are zero divisors.

\begin{examples}[Examples of Integral Domains]
    \label{ex:integral-domain}%
    \begin{itemize}
        \item Any field is an integral domain. Indeed, if \(a b = 0\) in a field, say \(b \neq 0\), then \(b\inverse\) exists, and
    \[
        a = a b b\inverse = a \cdot 0 = 0
    \]
    so \(a = 0\).
        \item Similarly, any subring of a field is also an integral domain. This is really how we should think of integral domains.
        \item The rings \(\ZZ\), \(\ZZ\bqty{i}\), \(\QQ\), \(\CC\), \(\RR\bqty{x}\) \(\ZZ\bqty{X}\) are integral domains.
        \item The rings \(\flatfrac{\ZZ}{6}\), \(\flatfrac{\ZZ}{\pqty{pq}}\), \(\flatfrac{\ZZ}{2} \times \flatfrac{\ZZ}{2}\), etc.
        \item In particular, direct products cannot be integral domains, since \(\pqty{0, 1}\) and \(\pqty{1, 0}\) product to \(0\).
    \end{itemize}
\end{examples}

\begin{lemma}[Finite Integral Domains are Finite Fields]
    \label{lem:finite-integral-domain-finate-field}%
    Let \(R\) be a finite integral domain. Then \(R\) is a field.
\end{lemma}

An example of this is \(\flatfrac{\ZZ}{p}\) for a prime \(p\).

\begin{proof}
    Let \(a \in R\) be non-zero, we want to prove it is a unit. Consider the multiplication function
    \[
        \function{\mu_a}{R}{R}{r}{ar.}
    \]

    This is a group homomorphism for the underlying group structure. Since \(R\) is an integral domain, by definition, we have
    \[
        \ker \pqty{\mu_a} = \Bqty{0}.
    \]

    So \(\mu_a\) is injective. Recall that \(R\) is finite, which means \(\mu_a\) is bijective.

    In particular, \(1 \in \img \pqty{\mu_a}\), so there is some \(b \in R\) with \(b \neq 0\), such that
    \[
        a \cdot b = 1
    \]
    so \(a\) is indeed a unit, finishing our proof.
\end{proof}

\begin{definition}[Field of Fractions]
    \label{def:field-of-fractions}%
    Let \(R\) be an integral domain. A \emph{field of fractions} of \(R\) is a field \(F\), such that
    \begin{enumerate}
        \item \(R \subring F\) is a subring, meaning there is an injective ring homomorphism from \(R\) to \(F\); and
        \item any \(x \in F\) can be written as \(a b \inverse\), for \(a, b \in R\), and \(b\inverse\) as the multiplicative inverse of \(b\) in \(F\).
    \end{enumerate}
\end{definition}

The first requirement is `sticking \(R\) into some field \(F\)', while the second requirement is `\(F\) is only contained of fractions in \(R\)'.

\begin{example*}[Field of Fractions of \(\ZZ\)]
    \label{ex:field-of-fractions-z}%
    The field \(\QQ\) is a field of fractions of \(\ZZ\).
\end{example*}

\begin{theorem}[Existence of Field of Fractions]
    \label{thm:existence-field-of-fractions}%
    Every integral domain \(R\) has a field of fractions.
\end{theorem}

Note \autoref{thm:existence-field-of-fractions} tells us that all integral rings are effectively subrings. The proof is inspired by the construction of \(\QQ\) by \(\ZZ\).

\begin{proof}
    Define a set
    \[
        S = \set{\pqty{a, b} \in R \times R}{b \neq 0}.
    \]

    We define a relation \(\sim\) over \(S\), being
    \[
        \pqty{a, b} \sim \pqty{c, d} \iff ad = bc,
    \]
    where here the multiplication is in \(R\).

    We want to show \(\sim\) is an equivalence relation. \emph{Symmetry} and \emph{reflective} are obvious. We conduct a non-trivial check on the \emph{transitivity}. Suppose that
    \[
        \pqty{a, b} \sim \pqty{c, d} \qand \pqty{c, d} \sim \pqty{e, f}.
    \]

    We have \(ad = bc\) and \(cf = de\). Multiplying the first equality by \(f\), and the second by \(b\), we have
    \[
        adf = bcf \qand bcf = bde,
    \]
    and hence \(adf = bde\). So we may rearrange and combine terms to get
    \[
        d \pqty{af - be} = 0.
    \]

    We have \(d \neq 0\), by our definition of \(S\). Using that \(R\) is an integral domain, we have \(af - be = 0\), so indeed \(af = be\), and we have \(\pqty{a, b} \sim \pqty{e, f}\).

    Now, we take the equivalence classes
    \[
        F = \flatfrac{S}{\sim},
    \]
    and we denote \(\frac{a}{b} = \bqty{\pqty{a, b}}\) for the conjugacy class. We make \(F\) a ring by the following operations.

    \[
        \frac{a}{b} + \frac{c}{d} = \frac{ad + bc}{bd} \qand \frac{a}{b} \cdot \frac{c}{d} = \frac{ac}{bd}.
    \]

    Elementary checks show that these operations are well-defined and \(F\) is a ring.

    To see \(F\) is a field, note that \(\frac{a}{b} \neq 0_F\), then \(\frac{a}{b} \neq \frac{0}{1}\), so \(a \cdot 1 \neq b \cdot 0\), so \(a \neq 0\). Then \(\frac{b}{a} \in F\) is an inverse of \(\frac{a}{b}\).

    Finally, we define a homomorphism
    \[
        \function{\phi}{R}{F}{r}{\frac{r}{1}.}
    \]

    Clearly, \(\phi\) is a ring homomorphism, and observe that
    \[
        \ker\pqty{\phi} = \set{r \in R}{\frac{r}{1} = 0} = \Bqty{0}.
    \]

    By \autoref{thm:first-isomorphism-theorem-ring} (First Isomorphism Theorem for Rings), we have \(\img \pqty{\phi} \isomorphic R\). Since every \(\frac{a}{b} \in R\) can be written as
    \[
        \frac{a}{b} = \frac{a}{1} \cdot \frac{1}{b} = \pqty{\frac{a}{1}} \cdot \pqty{\frac{b}{1}}\inverse, 
    \]
    so we have completed the proof.
\end{proof}

\begin{proposition}[Ideals of Fields]
    \label{prop:ideals-fields}%
    Let \(R\) be a ring. Then \(R\) is a field if and only if the only ideals of \(R\) are \(\Bqty{0}\) and \(R\).
\end{proposition}

\begin{proof}
    Let \(R\) be a field and take some ideal \(I \idealsub R\). Say \(I \neq \Bqty{0}\). Then \(I\) contains a unit, say \(u \in I\). But then \(1 \in I\) since \(1 = uv\) for some \(v \in R\). So \(I = R\).

    Conversely, take \(0 \neq r \in R\). Then we have \(\pqty{r} = R\), and hence \(1 \in R\) gives \(1 = r \cdot s\) for some \(s \in R\). So \(r\) is a unit, and \(R\) is a field.
\end{proof}

\begin{definition}[Maximal Ideal]
    \label{def:maximal-ideal}%
    An ideal \(I \idealsub R\) with \(I \neq R\) is \emph{maximal}, if for any \(J \idealsub R\), such that \(I \subseteq J \subseteq R\), we have either \(J = I\) or \(J = R\).
\end{definition}

Note that the zero ideal is maximal if and only if the ring itself is a field.

\begin{proposition}[Ideals are Maximal if and only if Quotient is Field]
    \label{prop:ideal-maximal-quotient-field}%
    An ideal \(I \idealsub R\) is maximal if and only if the quotient \(\flatfrac{R}{I}\) is a field.
\end{proposition}

\begin{proof}
    \(\flatfrac{R}{I}\) is a field if and only if its only ideals are the trivial ones, i.e., \(\Bqty{0}\) and \(\flatfrac{R}{I}\) itself, by \autoref{prop:ideals-fields}.

    Using the correspondence result for ideals between \(R\) and \(\flatfrac{R}{I}\) finishes the proof.
\end{proof}

\begin{definition}[Prime Ideal]
    \label{def:prime-ideal}%
    An ideal \(I \idealsub R\) with \(I \neq R\) is \emph{prime}, if whenever \(ab \in I\), we have \(a \in I\) or \(b \in I\).
\end{definition}

\begin{example}[Prime Ideals of \(\ZZ\)]
    \label{ex:prime-ideals-of-z}%
    An ideal \(n\ZZ \idealsub \ZZ\) is a prime ideal if and only if \(n\) is zero, or a prime.
    
    Indeed, if \(ab \in p\ZZ\) for some prime \(p\), then \(ab\) is a multiple of \(p\), so either \(a\) or \(b\) is also a multiple of \(p\).

    Conversely, if \(n\) is composite, consider \(n \ZZ\). Let \(n = n_1 n_2\) where \(1 < n_1, n_2 < n\). Then \(n_1\) and \(n_2\) product to a multiple of \(n\), in \(n \ZZ\), but neither are in the ideal \(n \ZZ\).
\end{example}

Note the similarity in this definition with \autoref{def:integral-domain-zero-divisor} (Integral Domain), i.e., a ring is an integral domain if and only if the zero ideal is prime.

\begin{proposition}[Ideals are Prime if and only if Quotient is Integral]
    \label{prop:ideal-prime-quotient-integral}%
    An ideal \(I \idealsub R\) is prime if and only if \(\flatfrac{R}{I}\) is an integral domain.
\end{proposition}

\begin{proof}
    If \(I\) is prime, let \(\pqty{a + I}\) and \(\pqty{b + I}\) be cosets in \(\flatfrac{R}{I}\) for some \(a, b \in R\). The product satisfies
    \[
        \pqty{a + I} \cdot \pqty{b + I} = ab + I = 0 + I
    \]
    is the zero ideal if and only if \(ab \in I\). This, since \(I\) is prime, implies that \(a \in I\) or \(b \in I\), so \(a + I = 0 + I\) or \(b + I = 0 + I\), showing that \(\flatfrac{R}{I}\) is indeed integral.

    Conversely, if \(\flatfrac{R}{I}\) is integral, consider \(ab \in I\) where \(ab = 0\). Then
    \[
        ab + I = \pqty{a + I} \pqty{b + I} = 0 + I
    \]
    in \(\flatfrac{R}{I}\) which is integral. Therefore, \(a + I = 0\) or \(b + I = 0\).
\end{proof}

\begin{proposition}[Maximal Ideals are Prime]
    \label{prop:maximal-ideal-prime}%
    Every maximal ideal is a prime ideal.
\end{proposition}

\begin{proof}
    Every field is an integral domain. \autoref{prop:ideal-maximal-quotient-field} and \autoref{prop:ideal-prime-quotient-integral} implies the result.
\end{proof}

\begin{remark}
    The converse is false. Consider the example \(\pqty{0} \idealsub \ZZ\), and \(\pqty{x} \idealsub \RR\bqty{x, y}\) are both prime but not maximal.
\end{remark}

\begin{proposition}[Characteristic of Integral Domain is Zero or Prime]
    \label{prop:characteristic-integral-domain}%
    The characteristic of an integral domain \(R\) is either \(0\) or prime.
\end{proposition}

\begin{proof}
    Consider the unique homomorphism \(\functiondomain{f}{{\ZZ}}{R}\). Characteristic \(n\) implies that \(\pqty{n} \idealsub \ZZ\) is the kernel.

    By \autoref{thm:first-isomorphism-theorem-ring} (First Isomorphism Theorem for Rings), we have \(\flatfrac{\ZZ}{n\ZZ} \subring R\). Then \(\flatfrac{\ZZ}{n\ZZ}\) is therefore integral, and hence \(\pqty{n}\) must be a prime ideal.
\end{proof}